% !PRINTED 27
\weberpage{27}

\begin{center}
\S.\,8.

\gesperrt{Metacyklische Gruppen.}
\end{center}

Wir haben in \S.\,176 des ersten Bandes bei den Permutationsgruppen
 metacyklische Gruppen definirt und ihre Bedeutung
für die Auflösung der Gleichungen kennen gelernt. Der Begriff
ist aber von der besonderen Bedeutung der Gruppenelemente
unabhängig, und wir definiren daher jetzt allgemein:

\gesperrt{Eine Gruppe, deren Indexreihe aus lauter Primzahlen
 besteht, soll eine metacyklische Gruppe heissen.}

Für diese Gruppen gilt ein für die Anwendungen auf die
Algebra wichtiger Satz, der aber von diesen Anwendungen unabhängig
 ist und den wir daher hier noch ableiten. Der Satz
lautet:

IV.\ \gesperrt{Eine metacyklische Gruppe hat einen von der
Einheitsgruppe verschiedenen commutativen
Normaltheiler}\footnote{C. J o r d a n, Traité des substitutions, Livre IV.}.

Ist $P$ eine metacyklische Gruppe und

\weq{1}{
\begin{array}{c}
P,\ P_1,\ P_2 \cdots P_{\mu-1},\ 1\\
j_1,\ j_2 \cdots j_{\mu-1},\ j_\mu
\end{array}
}

die Compositionsreihe, deren Indexreihe $j_1, j_2 \cdots$ aus lauter
Primzahlen besteht, so ist $P_{\mu-1}$ jedenfalls eine commutative
Gruppe; denn wenn $\pi$ irgend ein von 1 verschiedenes Element
aus $P_{\mu-1}$ ist, so sind die Potenzen

\weq{2}{1,\ \pi,\ \pi^2 \cdots \pi^{j_\mu - 1}}

(weil $j_\mu$ eine Primzahl ist) alle von einander verschieden und
machen also die ganze Gruppe $P_{\mu-1}$ aus; die cyklische Gruppe (2)
ist aber gewiss commutativ, weil für je zwei Exponenten $h, k$
immer

\[
\pi^h \pi^k = \pi^k \pi^h = \pi^{h+k}
\]

ist. Zugleich ist, nach der Definition der Compositionsreihe. % [sic: the full stop is as printed; a comma is intended.]
$P_{\mu-1}$ Normaltheiler von $P_{\mu-2}$.

Wir nehmen demnach jetzt an, wir haben einen von der
Einheit verschiedenen commutativen Normaltheiler $Q_\nu$ irgend
einer Gruppe $P_\nu$ der Reihe (1) und setzen ausserdem noch voraus,
%CONT>>
